\documentclass[a4paper]{article}

\usepackage{amssymb}
\usepackage{amsmath}
\usepackage{mathtools}
\usepackage{enumitem}


%opening
\title{Daniel's H3 Mathematics Notes}
\author{Daniel Ginting}
\date{\today}

\begin{document}
	
\maketitle

\newpage
\thispagestyle{empty}

\vspace*{\fill}
\begin{center}
	\emph{To poki.}
\end{center}
\vspace*{\fill}

\newpage

\begin{abstract}

This is my H3 Mathematics notes. This paper is nothing less than a re-formatting of Dr Jonathon Teo's H3 Slides in 2025. The contents therefore are not original to the author.

\end{abstract}




\section*{20\quad Nuclear Physics}

\subsection*{Content}

\begin{itemize}
    \item The nuclear atom
    \item Radioactive decay
    \item Nuclear processes and conservation laws
    \item Mass defect and nuclear binding energy
\end{itemize}

\subsection*{Learning Outcomes}

Candidates should be able to:

\begin{enumerate}[label=(\alph*)]
    \item infer from the results of the Rutherford $\alpha$-particle scattering experiment the existence and small size of the atomic nucleus
    \item distinguish between nucleon number (mass number) and proton number (atomic number)
    \item show an understanding that an element can exist in various isotopic forms, each with a different number of neutrons in the nucleus, and use the notation
    \[
        {}^{A}_{Z}X
    \]
    for the representation of nuclides
    \item show an understanding of the spontaneous and random nature of nuclear decay
    \item infer the random nature of radioactive decay from the fluctuations in count rate
    \item show an understanding of the origin and significance of background radiation
    \item show an understanding of the nature and properties of $\alpha$, $\beta$ and $\gamma$ radiations (knowledge of positron emission is not required)
    \item define the terms activity and decay constant and recall and solve problems using the equation
    \[
        A = \lambda N
    \]
    \item infer and sketch the exponential nature of radioactive decay and solve problems using the relationship
    \[
        x = x_0 e^{-\lambda t}
    \]
    where $x$ could represent activity, number of undecayed particles or received count rate
    \item define and use half-life as the time taken for a quantity $x$ to reduce to half its initial value
    \item solve problems using the relation
    \[
        \lambda = \frac{\ln 2}{t_{\frac{1}{2}}}
    \]
    \item discuss qualitatively the applications (e.g.\ medical and industrial uses) and hazards of radioactivity based on:
    \begin{enumerate}[label=(\roman*)]
        \item half-life of radioactive materials
        \item penetrating abilities and ionising effects of radioactive emissions
    \end{enumerate}
    \item represent simple nuclear reactions by nuclear equations of the form
    \[
        {}^{14}_{7}\mathrm{N} + {}^{4}_{2}\mathrm{He} \rightarrow {}^{17}_{8}\mathrm{O} + {}^{1}_{1}\mathrm{H}
    \]
    \item state and apply to problem solving the concept that nucleon number, charge and mass-energy are all conserved in nuclear processes
    \item show an understanding of how the conservation laws for energy and momentum in $\beta$ decay were used to predict the existence of the (anti)neutrino (knowledge of the antineutrino and the zoo of particles is not required)
    \item show an understanding of the concept of mass defect
    \item recall and apply the equivalence between energy and mass as represented by
    \[
        E = mc^2
    \]
    to solve problems
    \item show an understanding of the concept of nuclear binding energy and its relation to mass defect
    \item sketch the variation of binding energy per nucleon with nucleon number
    \item explain the relevance of binding energy per nucleon to nuclear fusion and to nuclear fission
\end{enumerate}






\end{document}
